% !TeX program = xelatex
% UTF-8; compile twice with XeLaTeX. No external images or data files needed.
\documentclass[a4paper,landscape,10pt]{article}
\usepackage[margin=14mm,headheight=14pt,headsep=5mm,footskip=7mm]{geometry}
\usepackage{fontspec}
\setmainfont{DejaVu Sans}
\newfontfamily\ipafont{DejaVu Sans}
\usepackage[table]{xcolor}
\usepackage{array,longtable,booktabs,fancyhdr}
\definecolor{ink}{HTML}{18384D}
\definecolor{band}{HTML}{E7EFF4}
\pagestyle{fancy}\fancyhf{}
\fancyhead[L]{\small\color{ink}Parlez Vous QSO?! · Nederlands / Dutch}
\fancyhead[R]{\small Englische Phrase · Nederlandse vertaling · IPA}
\fancyfoot[L]{\scriptsize Vorlage: Dutch v 1.0 · Credits: Shanondoa Bonifacio und Papi -- PE1MWO}
\fancyfoot[R]{\small\thepage}
\setlength{\parindent}{0pt}
\setlength{\tabcolsep}{3mm}
\renewcommand{\arraystretch}{1.25}
\newcommand{\ipa}[1]{{\ipafont[#1]}}
\begin{document}
{\LARGE\bfseries\color{ink}Parlez Vous QSO?!}\hfill{\large Nederlands / Dutch}

\vspace{2mm}
{\small QSO-Sprachführer mit niederländischer IPA-Lautschrift. Die Lautschrift bezieht sich auf die niederländische Spalte.}

\vspace{2mm}
{\footnotesize \textbf{Leseschlüssel:} \ipa{ˈ} Hauptbetonung, \ipa{ˌ} Nebenbetonung, \ipa{ː} langer Vokal;
\ipa{ɛi} ei/ij, \ipa{œy} ui, \ipa{ʌu} au/ou; \ipa{ɣ, x} g/ch, \ipa{ʋ} niederländisches w.
Grundlage: Standardniederländisch der Niederlande. Schrägstriche trennen Alternativen; \ipa{…} steht für einen Platzhalter.}

\vspace{1mm}
\fontsize{9}{11.8}\selectfont
\begin{longtable}{@{}>{\raggedright\arraybackslash}p{13mm}>{\raggedright\arraybackslash}p{73mm}>{\raggedright\arraybackslash}p{75mm}>{\raggedright\arraybackslash}p{90mm}@{}}
\toprule
\textbf{Ref.} & \textbf{English phrase} & \textbf{Nederlands} & \textbf{Lautschrift (IPA)}\\
\midrule
\endfirsthead
\toprule
\textbf{Ref.} & \textbf{English phrase} & \textbf{Nederlands} & \textbf{Lautschrift (IPA)}\\
\midrule
\endhead
\midrule
\multicolumn{4}{r}{\scriptsize Fortsetzung auf der nächsten Seite}\\
\endfoot
\bottomrule
\endlastfoot
\rowcolor{band} 1.0 & \textbf{Opening Greetings} & \multicolumn{2}{l}{\textbf{Openingszinnen}} \\*
1.1 & I am very pleased to meet you. & Aangenaam kennis te maken.\newline Leuk je te ontmoeten.\newline Ik vind het leuk je te ontmoeten. & \ipa{ˈaːŋɣənaːm ˈkɛnɪs tə ˈmaːkən / løːk jə tə ɔntˈmutən / ɪk vɪnt ət løːk jə tə ɔntˈmutən} \\
1.2 & Thank you for your call. & Bedankt voor je oproep. & \ipa{bəˈdɑŋkt voːr jə ˈɔprup} \\
1.3 & Good morning. & Goede morgen. & \ipa{ˈɣudə ˈmɔrɣən} \\
1.31 & Good day. & Goedendag & \ipa{ˌɣudənˈdɑx} \\
1.32 & Good afternoon. & Goede middag. & \ipa{ˈɣudə ˈmɪdɑx} \\
1.33 & Good evening. & Goede avond. & \ipa{ˈɣudə ˈaːvɔnt} \\
1.4 & I know only a little Dutch but I enjoy QSOs in Dutch. & Ik spreek maar een beetje Nederlands, maar ik vind het leuk te QSO’en in het Nederlands. & \ipa{ɪk spreːk maːr ən ˈbeːtjə ˈneːdərlɑnts maːr ɪk vɪnt ət løːk tə ky ɛs oːən ɪn ət ˈneːdərlɑnts} \\
\rowcolor{band} 2.0 & \textbf{Operator} & \multicolumn{2}{l}{\textbf{Operator}} \\*
2.1 & My name is Peter. & Mijn naam is Peter. & \ipa{mɛin naːm ɪs ˈpeːtər} \\
2.2 & I have been licensed since 1960. & Ik heb mijn vergunning sinds 1960. & \ipa{ɪk ɦɛp mɛin vərˈɣʏnɪŋ sɪnts ˈneːɣəntinɦɔndərtˌzɛstəx} \\
2.3 & I have been licensed for 43 years. & Ik heb mijn vergunning al 43 jaar. & \ipa{ɪk ɦɛp mɛin vərˈɣʏnɪŋ ɑl ˌdriːənˈveːrtəx jaːr} \\
2.4 & I live in Sittingbourne. & Ik woon in Sittingbourne. & \ipa{ɪk ʋoːn ɪn ˈsɪtɪŋbɔːn} \\
2.5 & Sittingbourne is a village. & Sittingbourne is een dorp. & \ipa{ˈsɪtɪŋbɔːn ɪs ən dɔrp} \\
2.51 & Sittingbourne is a town. & Sittingbourne is een stad. & \ipa{ˈsɪtɪŋbɔːn ɪs ən stɑt} \\
2.52 & Sittingbourne is a city. & Sittingbourne is een grote stad. & \ipa{ˈsɪtɪŋbɔːn ɪs ən ˈɣroːtə stɑt} \\
2.6 & Sittingbourne is between London and Dover. & Sittingbourne ligt tussen Londen en Dover. & \ipa{ˈsɪtɪŋbɔːn lɪxt ˈtʏsən ˈlɔndən ɛn ˈdəʊvə} \\
2.7 & I live in the South of England. & Ik woon in het Zuiden van Engeland. & \ipa{ɪk ʋoːn ɪn ət ˈzœydən vɑn ˈɛŋəlɑnt} \\
2.71 & I live in the North of England. & Ik woon in het Noorden van Engeland. & \ipa{ɪk ʋoːn ɪn ət ˈnoːrdən vɑn ˈɛŋəlɑnt} \\
2.72 & I live in the East of England. & Ik woon in het Oosten van Engeland. & \ipa{ɪk ʋoːn ɪn ət ˈoːstən vɑn ˈɛŋəlɑnt} \\
2.73 & I live in West of England. & Ik woon in het Westen van Engeland. & \ipa{ɪk ʋoːn ɪn ət ˈʋɛstən vɑn ˈɛŋəlɑnt} \\
2.74 & I live in South East England. & Ik woon in het Zuidoosten van Engeland. & \ipa{ɪk ʋoːn ɪn ət ˌzœytˈoːstən vɑn ˈɛŋəlɑnt} \\
2.75 & I live in South West England. & Ik woon in het Zuidwesten van Engeland. & \ipa{ɪk ʋoːn ɪn ət ˌzœytˈʋɛstən vɑn ˈɛŋəlɑnt} \\
2.76 & I live in North East England. & Ik woon in het Noordoosten van Engeland. & \ipa{ɪk ʋoːn ɪn ət ˌnoːrtˈoːstən vɑn ˈɛŋəlɑnt} \\
2.77 & I live in North West England. & Ik woon in het Noordwesten van Engeland. & \ipa{ɪk ʋoːn ɪn ət ˌnoːrtˈʋɛstən vɑn ˈɛŋəlɑnt} \\
2.78 & I live in central England. & Ik woon in het midden van Engeland. & \ipa{ɪk ʋoːn ɪn ət ˈmɪdən vɑn ˈɛŋəlɑnt} \\
2.8 & I live in the country. & Ik woon op het platteland. & \ipa{ɪk ʋoːn ɔp ət ˈplɑtəlɑnt} \\
2.81 & I live by the sea. & Ik woon aan de zee. & \ipa{ɪk ʋoːn aːn də zeː} \\
2.82 & I live near mountains. & Ik woon bij de bergen. & \ipa{ɪk ʋoːn bɛi də ˈbɛrɣən} \\
2.9 & Sittingbourne is 50Km from London. & Sittingbourne ligt 50 km van Londen. & \ipa{ˈsɪtɪŋbɔːn lɪxt ˈvɛiftəx ˌkiloːˈmeːtər vɑn ˈlɔndən} \\
2.11 & I am retired. & Ik ben gepensioeneerd.\newline Ik ben met pensioen. & \ipa{ɪk bɛn ɣəpɛnʃoːˈneːrt / ɪk bɛn mɛt pɛnˈʃun} \\
2.12 & My home call sign is M0ECS. & Mijn Call Sign thuis is M0ECS. & \ipa{mɛin kɔːl saɪn tœys ɪs ɛm nʏl eː seː ɛs} \\
\rowcolor{band} 3.0 & \textbf{Station} & \multicolumn{2}{l}{\textbf{Station}} \\*
3.1 & My radio is an Icom 706. & Mijn radiozender is een Icom 706.\newline Mijn radio is een Icom 706. & \ipa{mɛin ˈraːdijoːˌzɛndər ɪs ən ˈaɪkɒm ˌzeːvənɦɔndərtˈzɛs / mɛin ˈraːdijoː ɪs ən ˈaɪkɒm ˌzeːvənɦɔndərtˈzɛs} \\
3.2 & My station is all homebrew (homemade). & Mijn zender is helemaal zelfgemaakt.\newline Mijn station is zelfbouw. & \ipa{mɛin ˈzɛndər ɪs ˈɦeːləmaːl ˈzɛlfɣəmaːkt / mɛin staˈʃɔn ɪs ˈzɛlfbʌu} \\
3.3 & I am running 5 Watts. & Mijn vermogen is 5 Watt.\newline Ik gebruik 5 Watt. & \ipa{mɛin vərˈmoːɣən ɪs vɛif ʋɑt / ɪk ɣəˈbrœyk vɛif ʋɑt} \\
3.4 & My antenna is a long wire. & Mijn antenne is een draadantenne. & \ipa{mɛin ɑnˈtɛnə ɪs ən ˈdraːtɑnˌtɛnə} \\
3.41 & My antenna is a dipole. & Mijn antenne is een dipool. & \ipa{mɛin ɑnˈtɛnə ɪs ən diˈpoːl} \\
3.42 & My antenna is a G5RV. & Mijn antenne is een G5RV. & \ipa{mɛin ɑnˈtɛnə ɪs ən ɣeː vɛif ɛr veː} \\
3.43 & My antenna is a doublet. & Mijn antenne is een dubbel. & \ipa{mɛin ɑnˈtɛnə ɪs ən ˈdʏbəl} \\
3.44 & My antenna is a delta loop. & Mijn antenne is een delta loop. & \ipa{mɛin ɑnˈtɛnə ɪs ən ˈdɛltaː luːp} \\
3.45 & My antenna is a vertical. & Mijn antenne is een verticaal. & \ipa{mɛin ɑnˈtɛnə ɪs ən vɛrtiˈkaːl} \\
3.46 & My antenna is a log periodic. & Mijn antenne is een log periodiek. & \ipa{mɛin ɑnˈtɛnə ɪs ən lɔx peːriˈoːdik} \\
3.47 & My antenna is a Yagi beam. & Mijn antenne is een yagi beam. & \ipa{mɛin ɑnˈtɛnə ɪs ən ˈjaːɣi biːm} \\
3.48 & My antenna is a 5 element beam. & Mijn antenne is een 5 element beam. & \ipa{mɛin ɑnˈtɛnə ɪs ən vɛif eːləˈmɛnt biːm} \\
\rowcolor{band} 4.0 & \textbf{Conditions} & \multicolumn{2}{l}{\textbf{Condities}} \\*
4.1 & Thank you for my 599 report. & Bedankt voor mijn 599 rapport.\newline Dank u voor mijn 599 rapport. & \ipa{bəˈdɑŋkt voːr mɛin vɛif neːɣən neːɣən rɑˈpɔrt / dɑŋk y voːr mɛin vɛif neːɣən neːɣən rɑˈpɔrt} \\
4.2 & The band conditions are good. & De zendcondities zijn goed.\newline De band condities zijn goed. & \ipa{də ˈzɛntkɔnˌdisis zɛin ɣut / də bɑnt kɔnˈdisis zɛin ɣut} \\
4.3 & The band conditions are bad. & De zendcondities zijn slecht.\newline De band condities zijn slecht. & \ipa{də ˈzɛntkɔnˌdisis zɛin slɛxt / də bɑnt kɔnˈdisis zɛin slɛxt} \\
4.4 & There is strong interference here. & Er is sterke interferentie / storing hier.\newline Er is een sterke storing hier. & \ipa{ɛr ɪs ˈstɛrkə ˌɪntərfəˈrɛn(t)si / ˈstoːrɪŋ ɦiːr / ɛr ɪs ən ˈstɛrkə ˈstoːrɪŋ ɦiːr} \\
4.5 & Can you change to the 80m band? & Kun jij veranderen naar de 80m band?\newline Kun je overgaan naar 80m band? & \ipa{kʏn jɛi vərˈɑndərən naːr də ˈtɑxtəx meːtər bɑnt / kʏn jə ˈoːvərɣaːn naːr ˈtɑxtəx meːtər bɑnt} \\
4.6 & Please change frequency to 3.558 MHz. & Kun je de frequentie veranderen naar 3.558 MHz.\newline Alstublieft veranderen van frequentie naar 3.558 MHz. & \ipa{kʏn jə də freːˈkʋɛn(t)si vərˈɑndərən naːr driː ˈkɔmaː vɛif vɛif ɑxt ˈmeːɣaˌɦɛrts / ɑlstyˈblift vərˈɑndərən vɑn freːˈkʋɛn(t)si naːr driː ˈkɔmaː vɛif vɛif ɑxt ˈmeːɣaˌɦɛrts} \\
\rowcolor{band} 5.0 & \textbf{General Queries and Comments} & \multicolumn{2}{l}{\textbf{Algemene Vragen en Commentaren}} \\*
5.1 & What does <\#\#\#\#\#> mean please? & Wat betekent <\#\#\#\#\#> ? & \ipa{ʋɑt bəˈteːkənt …} \\
5.2 & Please wait a moment, I need to check my dictionary. & Wacht even a.u.b., ik kijk even in het woordenboek.\newline Wilt u even wachten, ik moet even in mijn woordenboek kijken. & \ipa{ʋɑxt ˈeːvən ˌɑlstyˈblift ɪk kɛik ˈeːvən ɪn ət ˈʋoːrdənˌbuk / ʋɪlt y ˈeːvən ˈʋɑxtən ɪk mut ˈeːvən ɪn mɛin ˈʋoːrdənˌbuk ˈkɛikən} \\
5.3 & How do you say, <\#\#\#\#\#> in Dutch? & Hoe zeg je <\#\#\#\#\#> in het Nederlands. & \ipa{ɦu zɛx jə … ɪn ət ˈneːdərlɑnts} \\
5.5 & My favourite modes of operation are CW / RTTY / PSK / MT63 / Olivia / Hellschreiber / Digital Modes / and SSB & Mijn favoriete gebruiksinstellingen zijn CW / RTTY / PSK / MT63 / Olivia / Hellschreiber / digitaal modes / en SSB. & \ipa{mɛin ˌfaːvoːˈritə ɣəˈbrœyksɪnˌstɛlɪŋən zɛin seː ˈdʏbəlveː / ɛr teː teː ɛi / peː ɛs kaː / ɛm teː ˌdriːənˈzɛstəx / oːˈliviˌaː / ˈhɛlˌʃʁaɪbɐ / diɣiˈtaːl ˈmoːdəs / ɛn ɛs ɛs beː} \\
5.6 & My best DX is Australia. & Mijn beste contact is Australië. & \ipa{mɛin ˈbɛstə kɔnˈtɑkt ɪs ʌusˈtraːliə} \\
5.7 & I enjoy working contests. & Ik houd van wedstrijden.\newline Ik vind het leuk om met wedstrijden mee te doen. & \ipa{ɪk ɦʌut vɑn ˈʋɛtstrɛidən / ɪk vɪnt ət løːk ɔm mɛt ˈʋɛtstrɛidən meː tə dun} \\
5.8 & I like working QRP. & Ik houd van werken met QRP. & \ipa{ɪk ɦʌut vɑn ˈʋɛrkən mɛt ky ɛr peː} \\
\rowcolor{band} 6.0 & \textbf{The Weather} & \multicolumn{2}{l}{\textbf{Het Weer}} \\*
6.1 & It is sunny. & Het is zonnig. & \ipa{ət ɪs ˈzɔnəx} \\
6.2 & It is windy. & Het waait. & \ipa{ət ʋaːit} \\
6.3 & It is raining. & Het regent. & \ipa{ət ˈreːɣənt} \\
6.4 & It is snowing. & Het sneeuwt. & \ipa{ət sneːʋt} \\
6.5 & It is very hot. & Het is erg warm. & \ipa{ət ɪs ɛrx ʋɑrm} \\
6.6 & It is very cold. & Het is erg koud. & \ipa{ət ɪs ɛrx kʌut} \\
6.7 & The temperature is 8C. & De temperatuur is 8 graden Celsius. & \ipa{də tɛmpəraːˈtyːr ɪs ɑxt ˈɣraːdən ˈsɛlsiʏs} \\
6.8 & It is overcast. & Het is bewolkt. & \ipa{ət ɪs bəˈʋɔlkt} \\
6.9 & It is foggy. & Het is mistig. & \ipa{ət ɪs ˈmɪstəx} \\
6.11 & It has been raining all day. & Het regent al de hele dag. & \ipa{ət ˈreːɣənt ɑl də ˈɦeːlə dɑx} \\
6.12 & There is 20cm of snow here. & Er ligt hier 20cm sneeuw. & \ipa{ɛr lɪxt ɦiːr ˈtʋɪntəx ˌsɛntiˈmeːtər sneːʋ} \\
\rowcolor{band} 8.0 & \textbf{QSL} & \multicolumn{2}{l}{\textbf{QSL}} \\*
8.1 & My QSL card is sure via the bureau. & Ik stuur mijn QSL kaart zeker via het bureau. & \ipa{ɪk styːr mɛin ky ɛs ɛl kaːrt ˈzeːkər ˈviːaː ət byˈroː} \\
8.2 & QSL via eQSL. & QSL via eQSL. & \ipa{ky ɛs ɛl ˈviːaː iː ky ɛs ɛl} \\
8.3 & No eQSL please. & Geen eQSL alstublieft. & \ipa{ɣeːn iː ky ɛs ɛl ɑlstyˈblift} \\
8.4 & I will send you my card direct. & Ik zal je mijn kaart direct sturen.\newline Ik stuur u mijn kaart direct. & \ipa{ɪk zɑl jə mɛin kaːrt diˈrɛkt ˈstyːrən / ɪk styːr y mɛin kaːrt diˈrɛkt} \\
8.5 & What is your address? & Wat is je adres?\newline Wat is uw adres? & \ipa{ʋɑt ɪs jə aːˈdrɛs / ʋɑt ɪs yʋ aːˈdrɛs} \\
8.6 & Please send your QSL card direct. & Kun je meteen je QSL kaart sturen, alsjeblieft.\newline Alstublieft stuur uw QSL kaart direct. & \ipa{kʏn jə məˈteːn jə ky ɛs ɛl kaːrt ˈstyːrən ɑlʃəˈblift / ɑlstyˈblift styːr yʋ ky ɛs ɛl kaːrt diˈrɛkt} \\
8.7 & My address is... & Mijn adres is: & \ipa{mɛin aːˈdrɛs ɪs …} \\
8.8 & My address is shown on QSL.com / EHam.net. & Mijn adres staat op QSL.com / Eham.net.\newline Mijn adres is te zien op QSL.com of op EHam.net. & \ipa{mɛin aːˈdrɛs staːt ɔp ky ɛs ɛl pʏnt kɔm / iː ɦɛm pʏnt nɛt / mɛin aːˈdrɛs ɪs tə zin ɔp ky ɛs ɛl pʏnt kɔm ɔf ɔp iː ɦɛm pʏnt nɛt} \\
8.9 & Sorry I do not send QSL cards. & Sorry, ik stuur geen QSL kaarten.\newline Sorry maar ik stuur geen QSL kaarten. & \ipa{ˈsɔri ɪk styːr ɣeːn ky ɛs ɛl ˈkaːrtən / ˈsɔri maːr ɪk styːr ɣeːn ky ɛs ɛl ˈkaːrtən} \\
8.11 & Please send QSL card via G1JQH. & Kunt u mij de QSL kaart via G1JQH sturen, alstublieft.\newline Alstublieft stuur de QSL kaart naar G1JQH. & \ipa{kʏnt y mɛi də ky ɛs ɛl kaːrt ˈviːaː ɣeː eːn jeː ky ɦaː ˈstyːrən ɑlstyˈblift / ɑlstyˈblift styːr də ky ɛs ɛl kaːrt naːr ɣeː eːn jeː ky ɦaː} \\
8.12 & My WACRAL / GQRP / DXCC / FISTS number is 1234. & Mijn WACRAL nummer is 1234.\newline Mijn WACRAL (registratienummer) van de groep / DXCC / FISTS nummer is 1234. & \ipa{mɛin ˈʋɑkraːl ˈnʏmər ɪs ˌdœyzənt ˌtʋeːɦɔndərt ˌviːrənˈdɛrtəx / mɛin ˈʋɑkraːl reːɣɪsˈtraːtsiˌnʏmər vɑn də ɣrup / deː ɪks seː seː / fɪsts ˈnʏmər ɪs ˌdœyzənt ˌtʋeːɦɔndərt ˌviːrənˈdɛrtəx} \\
\rowcolor{band} 9.0 & \textbf{Final Greetings} & \multicolumn{2}{l}{\textbf{Afsluitende groeten}} \\*
9.1 & Thank you. & Bedankt.\newline Dank u.\newline Dank u wel. & \ipa{bəˈdɑŋkt / dɑŋk y / dɑŋk y ʋɛl} \\
9.2 & Thank you for a very enjoyable contact. & Bedankt voor het aangename contact.\newline Dank je voor het plezierige contact. & \ipa{bəˈdɑŋkt voːr ət ˈaːŋɣənaːmə kɔnˈtɑkt / dɑŋk jə voːr ət pləˈzirəɣə kɔnˈtɑkt} \\
9.3 & I enjoyed our chat. & Ik vond het gesprek aangenaam / leuk / plezierig.\newline Ik vond het leuk om met je te chatten. & \ipa{ɪk vɔnt ət ɣəˈsprɛk ˈaːŋɣənaːm / løːk / pləˈzirəx / ɪk vɔnt ət løːk ɔm mɛt jə tə ˈtʃɛtən} \\
9.4 & I enjoyed our QSO in Dutch. & Ik vond het QSO in het Nederlands leuk.\newline Dank je wel voor je QSO in het Nederlands. & \ipa{ɪk vɔnt ət ky ɛs oː ɪn ət ˈneːdərlɑnts løːk / dɑŋk jə ʋɛl voːr jə ky ɛs oː ɪn ət ˈneːdərlɑnts} \\
9.5 & It has been great fun. & Het was erg leuk. & \ipa{ət ʋɑs ɛrx løːk} \\
9.6 & If you would like a regular contact with me, you can contact me via email. & Als je me vaker wilt spreken, kun je contact met me opnemen via e-mail.\newline Als je regelmatig contact met mij wilt, dan kun je mij bereiken via mijn e-mail. & \ipa{ɑls jə mə ˈvaːkər ʋɪlt ˈspreːkən kʏn jə kɔnˈtɑkt mɛt mə ˈɔpneːmən ˈviːaː ˈiːmeːl / ɑls jə ˈreːɣəlmaːtəx kɔnˈtɑkt mɛt mɛi ʋɪlt dɑn kʏn jə mɛi bəˈrɛikən ˈviːaː mɛin ˈiːmeːl} \\
9.7 & My email address is... & Mijn e-mail adres is… & \ipa{mɛin ˈiːmeːl aːˈdrɛs ɪs …} \\
9.8 & My email address is shown on QSL.com / EHam.net & Mijn e-mail adres kun je vinden op QSL.com / Eham.net.\newline Mijn e-mail adres is op de site QSL.com / EHam.net. & \ipa{mɛin ˈiːmeːl aːˈdrɛs kʏn jə ˈvɪndən ɔp ky ɛs ɛl pʏnt kɔm / iː ɦɛm pʏnt nɛt / mɛin ˈiːmeːl aːˈdrɛs ɪs ɔp də sɑit ky ɛs ɛl pʏnt kɔm / iː ɦɛm pʏnt nɛt} \\
9.9 & I look forward to meeting you again. & Ik hoop je snel weer te zien.\newline Ik verheug me je weer te ontmoeten. & \ipa{ɪk ɦoːp jə snɛl ʋeːr tə zin / ɪk vərˈɦøːx mə jə ʋeːr tə ɔntˈmutən} \\
9.11 & My best wishes to you and your family. & Ik wens jou en je familie het beste.\newline De groeten voor jou en je familie. & \ipa{ɪk ʋɛns jʌu ɛn jə faːˈmili ət ˈbɛstə / də ˈɣrutən voːr jʌu ɛn jə faːˈmili} \\
9.12 & Best wishes to your wife. & Groeten aan je vrouw. & \ipa{ˈɣrutən aːn jə vrʌu} \\
9.121 & Best wishes to your husband. & Groeten aan je man. & \ipa{ˈɣrutən aːn jə mɑn} \\
9.13 & Good luck. & Veel geluk.\newline Veel succes. & \ipa{veːl ɣəˈlʏk / veːl sʏkˈsɛs} \\
9.14 & Good health. & Gezondheid.\newline Goede gezondheid. & \ipa{ɣəˈzɔntɦɛit / ˈɣudə ɣəˈzɔntɦɛit} \\
9.15 & Goodbye. & Tot ziens. & \ipa{tɔt zins} \\
9.16 & I wish you an enjoyable afternoon. & Nog een prettige middag.\newline Ik wens je een prettige middag. & \ipa{nɔx ən ˈprɛtəɣə ˈmɪdɑx / ɪk ʋɛns jə ən ˈprɛtəɣə ˈmɪdɑx} \\
9.161 & I wish you an enjoyable evening. & Nog een prettige avond.\newline Ik wens je een prettige avond. & \ipa{nɔx ən ˈprɛtəɣə ˈaːvɔnt / ɪk ʋɛns jə ən ˈprɛtəɣə ˈaːvɔnt} \\
9.162 & I wish you an enjoyable weekend. & Nog een prettig weekend.\newline Ik wens je een prettig weekend. & \ipa{nɔx ən ˈprɛtəx ˈʋikɛnt / ɪk ʋɛns jə ən ˈprɛtəx ˈʋikɛnt} \\
9.17 & Enjoy your holiday. & Geniet van je vakantie. & \ipa{ɣəˈnit vɑn jə vaːˈkɑn(t)si} \\
9.18 & Happy New Year. & Gelukkig Nieuwjaar. & \ipa{ɣəˈlʏkəx ˈniʋjaːr} \\
\rowcolor{band} 10.0 & \textbf{Country} & \multicolumn{2}{l}{\textbf{Land}} \\*
10.1 & Albania & Albanië & \ipa{ɑlˈbaːniə} \\
10.2 & Austria & Oostenrijk & \ipa{ˈoːstənrɛik} \\
10.3 & Belgium & België & \ipa{ˈbɛlɣiə} \\
10.4 & Belorussia & Wit-Rusland & \ipa{ˈʋɪtˌrʏslɑnt} \\
10.5 & Bosnia \& Herzegovina & Bosnië Herzegovina & \ipa{ˈbɔsniə ɦɛrtseːˈɣoːvinaː} \\
10.6 & Bulgaria & Bulgarije & \ipa{bʏlˈɣaːriə} \\
10.7 & Croatia & Kroatië & \ipa{kroːˈaːtsiə} \\
10.8 & Czech Republic & Tsjechië & \ipa{ˈtʃɛxiə} \\
10.9 & Denmark & Denemarken & \ipa{ˈdeːnəˌmɑrkən} \\
10.11 & Eire & Ierland & \ipa{ˈiːrlɑnt} \\
10.12 & England & Engeland & \ipa{ˈɛŋəlɑnt} \\
10.13 & Estonia & Estonië & \ipa{ˈɛstoːniə} \\
10.14 & Finland & Finland & \ipa{ˈfɪnlɑnt} \\
10.15 & France & Frankrijk & \ipa{ˈfrɑŋkrɛik} \\
10.16 & Germany & Duitsland & \ipa{ˈdœytslɑnt} \\
10.17 & Gibraltar & Gibraltar & \ipa{ʒiˈbrɑltɑr} \\
10.18 & Greece & Griekenland & \ipa{ˈɣrikənlɑnt} \\
10.19 & Holland & Nederland & \ipa{ˈneːdərlɑnt} \\
10.21 & Hungary & Hongarije & \ipa{ɦɔŋˈɣaːriə} \\
10.22 & Iceland & IJsland & \ipa{ˈɛislɑnt} \\
10.23 & Northern Ireland & Noord-Ierland & \ipa{ˌnoːrtˈiːrlɑnt} \\
10.24 & Isle of Man & Isle of Man & \ipa{ˌaɪl əv ˈmæn} \\
10.25 & Italy & Italië & \ipa{iˈtaːliə} \\
10.26 & Latvia & Letland & \ipa{ˈlɛtlɑnt} \\
10.27 & Liechtenstein & Liechtenstein & \ipa{ˈlɪçtn̩ʃtaɪn} \\
10.28 & Lithuania & Litouwen & \ipa{ˈlitʌuən} \\
10.29 & Luxembourg & Luxemburg & \ipa{ˈlʏksəmˌbʏrx} \\
10.31 & Macedonia & Macedonië & \ipa{maːseːˈdoːniə} \\
10.32 & Moldavia & Moldavië & \ipa{mɔlˈdaːviə} \\
10.33 & Netherlands & Nederland & \ipa{ˈneːdərlɑnt} \\
10.34 & Norway & Noorwegen & \ipa{ˈnoːrʋeːɣən} \\
10.35 & Poland & Polen & \ipa{ˈpoːlən} \\
10.36 & Portugal & Portugal & \ipa{ˈpɔrtyɣɑl} \\
10.37 & Romania & Roemenië & \ipa{ruˈmeːniə} \\
10.38 & Russian Federation & Rusland & \ipa{ˈrʏslɑnt} \\
10.39 & Scotland & Schotland & \ipa{ˈsxɔtlɑnt} \\
10.41 & Slovakia & Slowakije & \ipa{sloːˈʋaːkiə} \\
10.42 & Slovenia & Slovenië & \ipa{sloːˈveːniə} \\
10.43 & Spain & Spanje & \ipa{ˈspɑnjə} \\
10.44 & Svalbard & Svalbard & \ipa{ˈsvɑlbɑrt} \\
10.45 & Sweden & Zweden & \ipa{ˈzʋeːdən} \\
10.46 & Switzerland & Zwitserland & \ipa{ˈzʋɪtsərlɑnt} \\
10.47 & Turkey & Turkije & \ipa{ˈtʏrkiə} \\
10.48 & Ukraine & Oekraïne & \ipa{ukraːˈinə} \\
10.49 & Wales & Wales & \ipa{weɪlz} \\
\end{longtable}
\normalsize
\section*{Hinweise zur Übertragung und Aussprache}
\fontsize{8.5}{10.5}\selectfont\raggedright
\textbf{Vorlage und Umfang.} Grundlage ist „Dutch v 1.0.doc“. Credits der Quelle: Shanondoa Bonifacio
(ohne Rufzeichen) und Papi, PE1MWO. Alle 142 Phrasen und Ländernamen sowie die unterschiedlichen
Formulierungsvarianten sind enthalten. Nummernsprünge und der fehlende Abschnitt 7 bleiben erhalten.

\smallskip
\textbf{Dokumentierte Korrekturen.} „pension“ wurde zu „pensioen“ (2.11), „alsublieft“ zu „alstublieft“ (8.11),
„WALCRAL“ zu „WACRAL“ (8.12), „contakt“ zu „contact“ (9.6), „Mij e-mail“ zu „Mijn e-mail“ und „side“ zu „site“ (9.8),
„success“ zu „succes“ (9.13) sowie „Ukraïne“ zu „Oekraïne“ (10.48). In 10.2 wurde das falsche „Australië“
bei Austria durch „Oostenrijk“ ersetzt. Die drei mit einer Tilde markierten Lücken wurden ergänzt:
Wit-Rusland (10.4), Isle of Man (10.24), Svalbard (10.44). Australien in 5.6 bleibt korrekt als „Australië“ stehen.
In 8.2 wurde das fehlende e in eQSL ergänzt; 8.3 wurde gemäß „No eQSL please“ als „Geen eQSL alstublieft“ berichtigt.
Die kaufmännischen Und-Zeichen zwischen alternativen Sätzen wurden durch Zeilenumbrüche ersetzt,
Mehrfachleerzeichen entfernt und die Webadressen in 8.8 einheitlich geschrieben.

\smallskip
\textbf{Inhaltliche Eigenheiten.} Andere ungewöhnliche Formulierungen bleiben erkennbar erhalten,
etwa „dubbel“ für Doublet (3.43), „digitaal modes“ (5.5) und die unvollständige zweite Variante in 8.12.
Dort fehlt GQRP in der niederländischen Quelle; es wurde nicht stillschweigend eingesetzt.
Die Betriebsartenbezeichnung und die unterschiedliche Bedeutung mancher Alternativsätze wurden nicht umgeschrieben.
Historische Ländernamen und Webadressen entsprechen weiterhin der Vorlage. Holland und Netherlands stehen beide als „Nederland“.

\smallskip
\textbf{Aussprachebasis.} Die IPA-Spalte ist eine breite, manuell erstellte Aussprachehilfe für das
Standardniederländische der Niederlande, keine enge Transkription einer Aufnahme. \ipa{ˈ} markiert Haupt-,
\ipa{ˌ} ausgewählte Nebenbetonungen. Bei \ipa{eː, oː, øː} ist eine einfache lange Vokalschreibung gewählt;
in vielen niederländischen Aussprachen werden diese Vokale leicht diphthongiert.
\ipa{ɛi, œy, ʌu} sind die Diphthonge von ei/ij, ui und au/ou. Kurze \ipa{ɪ, ʏ, ɑ} unterscheiden sich von
\ipa{i, y, aː}; i, u und y sind nicht automatisch lang. Vor r sind hier häufig längere Varianten angegeben.

\smallskip
\textbf{Konsonanten und Reduktionen.} g ist überwiegend als \ipa{ɣ}, ch beziehungsweise auslautendes g als
\ipa{x} dargestellt; nördliche Sprecher können auch \ipa{ɣ} stimmlos und weiter hinten sprechen.
Das w steht als \ipa{ʋ}, r als vereinfachtes \ipa{r}; regionale r-Varianten und flämische Lautvarianten sind nicht einzeln markiert.
Auslautverhärtung ist berücksichtigt. Unbetonte Wörter stehen häufig reduziert: „het“ \ipa{ət}, „een“ \ipa{ən},
„je“ \ipa{jə}, „me“ \ipa{mə}. Das n in Endungen auf -en ist zur Orientierung meist beibehalten,
obwohl es im Gespräch oft entfällt. Die Klammer in \ipa{n(t)si} kennzeichnet ein optionales t.
Assimilation und weitere Reduktionen im zusammenhängenden Sprechen sind nicht vollständig dargestellt.

\smallskip
\textbf{Namen, Rufzeichen und Fachwörter.} Sittingbourne, Dover, Isle of Man, Wales, Call Sign,
Icom, Beam und Loop sind mit einer gebräuchlichen englischen Aussprache dargestellt; Hellschreiber und
Liechtenstein mit deutscher Aussprache. Londen, Engeland und die meisten anderen Namen werden niederländisch gelesen.
Rufzeichen und die meisten Kürzel stehen mit niederländischen Buchstabennamen: G „gee“, J „jee“, Q „kuu“, W „dubbel-wee“, Y „ij“;
das internationale Buchstabieralphabet ist nicht ausgeschrieben. WACRAL und FISTS werden als Wörter gelesen.
QSO’en steht für die Buchstaben Q S O plus die Endung -en. Das e in eQSL und eHam wird als \ipa{iː} gelesen;
Webpunkte als „punt“, com und net als Wörter. „a.u.b.“ ist als „alstublieft“ ausgeschrieben.

\smallskip
\textbf{Zahlen und Platzhalter.} 1960 wird als „negentienhonderdzestig“, 43 als „drieënveertig“,
706 als „zevenhonderdzes“, MT63 mit „drieënzestig“, 599 als „vijf negen negen“ und
1234 als „duizend tweehonderd vierendertig“ gelesen. 3.558 MHz ist im QSO-Kontext „drie komma vijf vijf acht megahertz“;
80 m, 50 km und 20 cm stehen mit „tachtig meter“, „vijftig kilometer“ und „twintig centimeter“.
Die Platzhalter \texttt{<\#\#\#\#\#>} und offene Adressangaben erscheinen in der Lautschrift als Auslassung.

\smallskip
\textbf{Fachhinweise.} Carlos Gussenhoven: „Dutch“, Journal of the International Phonetic Association
22(1--2), 1992, S. 45--47, DOI 10.1017/S002510030000459X (bibliografischer Hinweis zur IPA-Beschreibung).
Ländernamen: Nederlandse Taalunie, „Buitenlandse Aardrijkskundige Namen“, namen.taalunie.org.
Die einzelnen Phrasen wurden manuell transkribiert; bei Eigennamen, Marken und Kürzeln sind Varianten möglich.
\end{document}
